\subsection{Stopping and starting deliberately}
\label{sec:ha:disabled:stopping}

Not every stop is a failure. An operator takes a component down for a release,
for a configuration change, or because the trading day has ended. None of those
should cost the venue what a crash costs.

\begin{requirement}{R-0034}{An orderly stop leaves less to recover than a crash}
  A component asked to stop shall record its state before it goes, so that what
  it has to recover afterwards is no more than what it wrote.
  \rationale{A planned restart is the common case, for a release or a configuration change,
  and it need not cost what an unplanned restart does. It is also the cheapest way
  to exercise the recovery path regularly.}
  \uncovered{no scenario stops a component deliberately, and nothing would be written if one did}
\end{requirement}

\begin{requirement}{R-0035}{The start and end of a trading day are published}
  The venue shall publish the start of a trading day and its end as events that
  any interested party may subscribe to.
  \rationale{A venue that has just started is either resuming a trading day or beginning a new
  one, and which of those it is changes what every recovery in this section means.
  A component must be able to ask rather than infer it. Publishing the event once
  to every subscriber also stops each component being told separately, which is how
  components come to hold different answers.}
  \uncovered{no scenario spans a trading day boundary, which the venue does not mark}
\end{requirement}
